% !TeX program = lualatex
% Standalone notebook; compile twice: lualatex 38.tex
% Research batch 39 down to 25, 27 September 2026.
% Original section preserved verbatim except for marked insertion blocks.
\documentclass[11pt,a4paper]{article}
\usepackage[margin=24mm,headheight=14pt]{geometry}
\usepackage{fontspec}
\setmainfont{Latin Modern Roman}
\setsansfont{Latin Modern Sans}
\setmonofont{Latin Modern Mono}
\usepackage{amsmath,amssymb,mathtools}
\usepackage{unicode-math}
\setmathfont{Latin Modern Math}
\usepackage{microtype}
\usepackage{enumitem,xurl,needspace}
\usepackage[dvipsnames]{xcolor}
\usepackage{hyperref}
\hypersetup{unicode=true,colorlinks=true,linkcolor=MidnightBlue,urlcolor=MidnightBlue,
 pdftitle={Research notebook: Problem 38},
 pdfauthor={Open Problems Math research notebook},bookmarksnumbered=true}
\usepackage{bookmark,titlesec,fancyhdr}
\titleformat{\section}{\Large\bfseries\raggedright}{\thesection}{0.7em}{}
\pagestyle{fancy}\fancyhf{}
\fancyhead[L]{\small\sffamily Open Problems Math\enspace|\enspace Research notebook}
\fancyhead[R]{\small\sffamily Problem 38}
\fancyfoot[L]{\small\sffamily 27 September 2026}
\fancyfoot[R]{\small\sffamily \thepage}
\setlength{\parindent}{0pt}
\setlength{\parskip}{5pt plus 1pt}
\setlength{\emergencystretch}{3em}
\setcounter{secnumdepth}{1}
\setcounter{tocdepth}{1}
\setlist[itemize]{leftmargin=3em,itemsep=5pt,topsep=4pt,parsep=0pt}
\allowdisplaybreaks
\newcommand{\N}{\mathbb{N}}
\newcommand{\Z}{\mathbb{Z}}
\newcommand{\Q}{\mathbb{Q}}
\newcommand{\R}{\mathbb{R}}
\newcommand{\C}{\mathbb{C}}
\newcommand{\F}{\mathbb{F}}
\newcommand{\A}{\mathbb{A}}
\newcommand{\PP}{\mathbb P}
\newcommand{\E}{\mathbb{E}}
\newcommand{\eps}{\varepsilon}
\newcommand{\dd}{\,\mathrm d}
\newcommand{\sref}[2]{\hyperlink{source:#1:#2}{[S#2]}}
\newcommand{\eref}[2]{\hyperlink{extra:#1:#2}{[E#2]}}
\newif\ifshowcatalogue
\showcataloguetrue
\newcommand{\cataloguescope}[1]{\ifshowcatalogue
 \paragraph{Exact scope in the supplied catalogue.}
 \begin{quote}\small #1\end{quote}\fi}
\numberwithin{equation}{section}
\begin{document}
\setcounter{section}{37}
% BEGIN PRESERVED SECTION
% ================================================================
% problemId: problem.higher-dimensional-euclidean-kakeya-conjecture
% source releaseRank: 38; source releaseStatus: open_with_solved_subcases
% exactTarget SHA-256: 34c6a5eaecbaddea762aabc58019d3653906198e5c52d9123a6d1a03c2a566de
\clearpage
\section{\texorpdfstring{Higher-dimensional Euclidean Kakeya conjecture}{Higher-dimensional Euclidean Kakeya conjecture}}\label{38}
\subsection{Definitions and mathematical statement}
For $s\ge0$, Hausdorff measure is obtained from covers by sets of diameter tending to zero, weighting each covering set by its diameter to power $s$. Hausdorff dimension is $\dim_H E=\inf\{s:\mathcal H^s(E)=0\}$. A compact Kakeya set satisfies
\[
 \forall v\in S^{n-1}\ \exists a_v\in\R^n:\quad
 \{a_v+tv:0\le t\le1\}\subseteq E.
\]
The selected assertion is $\dim_H E=n$ for every such $E$ and every integer $n\ge4$. Positive Lebesgue measure is not required, and replacing Hausdorff dimension by a different dimension would change the statement.
\par\smallskip\textit{Formulation and concept references:} \sref{38}{1}.
\cataloguescope{Prove or disprove that for every integer n >= 4, every compact set E in R\textasciicircum{}n that contains a unit line segment in every direction has Hausdorff dimension n.}
\subsection{Short English statement}
Must a compact set containing a unit segment in every direction have full Hausdorff dimension in every dimension at least four?
% BEGIN RESEARCH 38
\subsection{Research attempt}

\paragraph{Current frontier.}
Wang--Zahl prove the three-dimensional set conjecture; Guth's 2026 Bourbaki
survey explains that proof. This is not the $n\ge4$ assertion selected here,
nor is the set theorem the full Kakeya maximal-function estimate.
\eref{38}{1}, \eref{38}{2}. The formulation source proves Hausdorff dimension
at least $3.25$ for \emph{sticky} Kakeya sets in $\R^4$, and describes the
distinct general-set bound obtained by the planebrush method. Stickiness
cannot be inserted as a hypothesis of this entry. \sref{38}{1}.
The September 2026 paper on lifting curved sets concerns a transfer of
geometric problems, not a theorem giving dimension $n$ for every linear
Kakeya set in every dimension. Its abstract was checked. \sref{38}{3}.

\paragraph{Attack route.}
Projection to three dimensions preserves some line segments but only gives
a three-dimensional lower bound, not the missing fiber dimension. Pairwise
tube intersections are directly computable but appear to lose too much.
The chosen route combines a scale-uniform directional estimate with a
weighted allocation of directions to scales. The point of the weighting is
to reach Hausdorff dimension, where covers have unequal radii, rather than
silently prove only a Minkowski statement.

\paragraph{Attempt.}
For $0<\delta<1$, let $T_\delta(a,v)$ be the radius-$\delta$ cylinder around
$\{a+tv:0\le t\le1\}$ and define
\[
 K_\delta f(v)=\sup_{a\in\R^n}\frac1{|T_\delta(a,v)|}
                      \int_{T_\delta(a,v)}|f(x)|\,dx.
\]
The following \emph{unproved hypothesis} is deliberately stronger than the
set statement:
\begin{equation}\label{r38:maximal}
 \|K_\delta f\|_{L^n(S^{n-1})}
 \le C_{n,\eta}\delta^{-\eta}\|f\|_{L^n(\R^n)}
 \quad\text{for every }\eta>0\text{ and every }\delta,f.
\end{equation}
It is a standard maximal-function target in the Kakeya framework discussed
in \eref{38}{2}. What follows is an explicit deduction from it, not a proof
of it.

\textbf{Lemma 38.1 (unequal-scale transfer).}
For a fixed $n$, \eqref{r38:maximal} implies the exact Hausdorff-dimension
conclusion of this entry in that dimension.

\emph{Proof.}
Fix $s<n$ and choose $\eta>0$ with $n-n\eta>s$. Cover a compact Kakeya set
$E$ by balls $B_i$ with radii $r_i\le2^{-J}$, grouping them so that
$2^{-j-1}<r_i\le2^{-j}=\delta_j$. Write $N_j$ for the number in group $j$,
$U_j$ for their union, and $U'_j$ for the union with each ball replaced by
the concentric ball of radius $3\delta_j$. Covers with infinite $s$-cost
need no argument; otherwise every $N_j$ is finite.

For each direction $v$, select only for this argument one unit segment
$\ell_v\subset E$. Since the cover contains every point of $\ell_v$,
\[
 \sum_{j\ge J}\mathcal H^1(\ell_v\cap U_j)\ge1.
\]
Choose fixed positive weights $w_j=c/(j+1)^2$, $j\ge0$, with
$\sum_{j\ge0}w_j=1/2$. Some $j\ge J$ has
$\mathcal H^1(\ell_v\cap U_j)\ge w_j$. Thickening the corresponding points
of the line in perpendicular disks of radius $\delta_j$ remains inside
$U'_j$, so $K_{\delta_j}1_{U'_j}(v)\ge w_j$. This uses Fubini in cylindrical
coordinates; endpoints have zero one-dimensional measure.
Define the measurable set
$D_j=\{v:K_{\delta_j}1_{U'_j}(v)\ge w_j\}$. No measurable choice of
$\ell_v$ is needed: these maximal-function level sets contain all the chosen
directions. Measurability follows by taking the supremum over a countable
dense set of translations; integrals of bounded indicators over translated
cylinders are continuous in the translation and direction.

The $D_j$ cover the sphere. Chebyshev's inequality and
\eqref{r38:maximal} yield
\[
 |S^{n-1}|\le\sum_{j\ge J}|D_j|
 \le C\sum_{j\ge J}(j+1)^{2n}\delta_j^{-n\eta}|U'_j|
 \le C'\sum_{j\ge J}(j+1)^{2n}N_j\delta_j^{n-n\eta}.
\]
Set $\gamma=n-n\eta-s>0$. Since
$\sup_{j\ge J}(j+1)^{2n}2^{-\gamma j}\to0$, this inequality gives a
positive lower bound, independent of the cover, on
$\sum_jN_j\delta_j^s$. Within a group $\delta_j<2r_i$, so the same holds
for the original Hausdorff cost up to a fixed factor. In particular
$\mathcal H^s(E)>0$, indeed the bound increases as $J\to\infty$.
Thus $\dim_HE\ge s$ for all $s<n$, proving $\dim_HE=n$.
\hfill$\square$

Can a simple overlap estimate prove the missing hypothesis? Take one unit
$\delta$-tube in each direction of a $\delta$-separated net of unoriented directions,
with $M\asymp\delta^{-(n-1)}$ tubes, and set $F=\sum_T1_T$. Elementary
cylinder geometry gives
\[
 |T\cap T'|\le C_n\frac{\delta^n}{\theta(T,T')}
 \quad(\theta\ge\delta).
\]
Here $\theta$ is the acute angle between unoriented axes. For one tube,
the number of directions at angular distance at most $\theta$
is $O((\theta/\delta)^{n-1})$. Summing over dyadic angular shells therefore
gives, for $n>2$,
\[
 \sum_{T'}|T\cap T'|\le C_n\delta,
 \qquad \int F^2\le C_nM\delta\asymp\delta^{-(n-2)}.
\]
The diagonal term $|T|\asymp\delta^{n-1}$ is harmless. Since
$\int F\asymp1$, Cauchy--Schwarz gives only
\[
 \left|\bigcup T\right|\gtrsim\delta^{n-2}.
\]
This basic calculation reaches a two-dimensional scale exponent, not $n$.
It is not a new best lower bound; it diagnoses the loss in the particular
pairwise route. Full ambient dimension would require bounds arbitrarily
close to $\delta^0$, with a way to control heavily shaded subsets across
scales. The lost information concerns simultaneous configurations of many
tubes, exactly the type of structure not recorded by pairwise angles alone.

\paragraph{Stress test.}
A lower bound for $|E_\delta|$ at equal scales need not establish Hausdorff
dimension. Lemma~38.1 explicitly handles arbitrary countable ball covers.
The constant in \eqref{r38:maximal} must be independent of tube locations,
$\delta$, and the chosen cover; making it depend on the set at a shrinking
scale defeats the proof. A family of tubes restricted to a hyperplane is
not a counterexample to the target because it lacks most directions.
Conversely, projecting a genuine Kakeya set to a three-plane provides no
positive lower bound on the dimensions of its fibers. Such a bound cannot
be inferred by subtracting dimensions without an appropriate slicing theorem
and a verified fiber hypothesis.

The strong maximal estimate is not claimed as a necessary condition for the
set theorem, so failing to prove it does not disprove Kakeya. The elementary
overlap argument provides neither a candidate exceptional compact set nor
the multiscale high-order estimate needed to replace \eqref{r38:maximal}.

\Needspace{8\baselineskip}
\paragraph{Outcome.}
\textbf{Outcome: Conditional reduction.}

The all-scale maximal hypothesis is shown to imply the exact Hausdorff target,
with the unequal-radius and direction-allocation details supplied. This is a
worked-out version of a classical implication, not a novelty claim. The
attempt also locates the exponent loss of its elementary two-tube strategy.
No higher-dimensional maximal estimate or counterexample set is produced;
that remains the substantive unresolved step.

% END RESEARCH 38

\subsection{Sources}
\begin{itemize}\raggedright
\item[\textnormal{[S1]}]\hypertarget{source:38:1}{} Mukul Rai Choudhuri, Revista Matematica Iberoamericana 42 (2026), 1227-1256.
\url{https://ems.press/journals/rmi/articles/14299708}.
{\small\textit{Catalogue formulation source.}}
\item[\textnormal{[S2]}]\hypertarget{source:38:2}{} Fourier analytic properties of Kakeya sets in finite fields.
\url{https://londmathsoc.onlinelibrary.wiley.com/doi/10.1112/blms.70367}.
{\small\textit{Catalogue research/status reference; not a proof endorsement.}}
\item[\textnormal{[S3]}]\hypertarget{source:38:3}{} Lifting curved Kakeya sets to linear Kakeya sets.
\url{https://arxiv.org/abs/2609.06299}.
{\small\textit{Catalogue research/status reference; not a proof endorsement.}}
\item[\textnormal{[S4]}]\hypertarget{source:38:4}{} A Unified Geometric Algebra Framework for the Kakeya Conjecture in All Dimensions and Links to the Riemann Zeta Function.
\url{https://www.preprints.org/manuscript/202602.1025}.
{\small\textit{Catalogue research/status reference; not a proof endorsement.}}
% BEGIN ADDED REFERENCES 38
\item[\textnormal{[E1]}]\hypertarget{extra:38:1}{} Hong Wang and Joshua Zahl,
\emph{Volume estimates for unions of convex sets, and the Kakeya set conjecture in three dimensions}, arXiv:2502.17655, main theorem.
\url{https://arxiv.org/abs/2502.17655}.
\item[\textnormal{[E2]}]\hypertarget{extra:38:2}{} Larry Guth,
\emph{The Kakeya conjecture, after Wang and Zahl}, arXiv:2604.03416,
Bourbaki survey, introduction and discussion of dimensional and maximal formulations.
\url{https://arxiv.org/abs/2604.03416}.
{\small\textit{References consulted 27 September 2026; no extension from dimension three to the present range is assumed.}}

% END ADDED REFERENCES 38
\end{itemize}

% END PRESERVED SECTION
\end{document}
